%======================================================================
\chapter{Results and Analysis}
%======================================================================

\subsection*{Baseline Infrastructure Analysis}

The current charging network in the \glsxtrshort{kwc-cma} region, serving $637,730$ residents \cite{statcan_census_2021} across $1,382.1 \text{ km}^2$ \cite{statcan_census_2021}, demonstrates significant coverage limitations. Our analysis revealed that the existing 183 charging stations \cite{ocm_2024} provide inadequate coverage, with \gls{l2} charging accessible to only 14.57\% of the population and \gls{l3} charging reaching a mere 1.78\%. This baseline assessment provided crucial context for optimization targets and infrastructure enhancement strategies.

\subsection*{Optimization Results}

The \glsxtrshort{mip} model generated a comprehensive solution addressing current coverage deficiencies while maintaining cost efficiency. The optimized network design incorporates 44 new \gls{l2} charging stations with single-port configurations in residential areas and community hubs, complemented by 11 new \gls{l3} fast-charging stations with dual-port configurations along major transportation corridors. Additionally, the model identified four existing \gls{l2} stations for strategic upgrade to \gls{l3} capability, augmenting the fast-charging network with eight additional high-capacity ports.

A key feature of our solution is the port retention strategy during upgrades, which ensures continued \gls{l2} charging availability by maintaining a minimum of one \gls{l2} port at each upgraded location. This approach balances the benefits of enhanced charging capabilities with the need to serve vehicles that are not compatible with \gls{l3} charging, while optimizing space utilization and implementation costs.

Implementation of the optimized solution would dramatically improve charging accessibility. \gls{l2} coverage would increase from \fbox{14.57\% to 96.14\%} of the population, while \gls{l3} coverage would expand from \fbox{1.78\% to 85.40\%}. Total charging capacity would increase to 463 \gls{l2} ports and 96 \gls{l3} ports, providing robust infrastructure for anticipated \glsxtrshort{ev} adoption growth.

\subsection*{Financial Analysis}

Economic analysis of the optimized solution revealed a carefully balanced infrastructure investment approach. The total infrastructure investment requirement of \$4,100,000 encompasses \$1,100,000 for new \gls{l2} stations, \$2,200,000 for new \gls{l3} stations, and \$800,000 for \gls{l2} to \gls{l3} upgrades. Equipment resale opportunities of \$105,000 reduce the net investment to \fbox{\$3,995,000}. 

This cost structure, detailed in Figure \ref{fig:2-optimization-results-plot} of the Appendix, demonstrates efficient resource allocation while achieving substantial coverage improvements.

\subsection*{Theoretical Coverage and Sensitivity Analysis}

The model revealed theoretical maximum coverage potentials of 95.51\% for \gls{l2} and 98.42\% for \gls{l3} accessibility. Our optimization achieves near-optimal \gls{l2} coverage at 96.14\% and substantial \gls{l3} coverage at 85.40\% while respecting practical constraints. 

Sensitivity analysis demonstrated robust solution stability with budget utilization at 99.88\% and coverage metrics exceeding minimum targets. See Figure \ref{fig:1-sensitivity-plot} of the Appendix.

\subsection*{Implementation Strategy}

The optimization results informed a three-phase implementation strategy designed to systematically enhance the charging network while minimizing operational disruptions. 

Phase 1, spanning months 1-4, focuses on strategic \gls{l2} to \gls{l3} upgrades, prioritizing grid capacity enhancements and equipment installation at identified high-impact locations. Phase 2 (months 5-8) encompasses new \gls{l3} station deployments, while Phase 3 (months 9-12) completes the \gls{l2} network expansion. 

This phased approach, illustrated in Figure \ref{fig:implementation-plan-plot} of the Appendix, ensures coordinated infrastructure development and efficient resource utilization.